\documentclass[11pt]{article}

% parameters are fairly self-explanatory here, thankfully
\usepackage[width=7in,height=9.5in,top=0.75in,centering]{geometry}

% for math-ing
\usepackage{amsmath}
\usepackage{amssymb}


%%%%%% tables %%%%%%%%

\usepackage{array}
\renewcommand{\arraystretch}{1.5}
\setlength{\tabcolsep}{8pt}

%%%%%% plotting %%%%%%

\usepackage[dvipsnames]{xcolor}
\usepackage{tikz}
\usepackage{pgfplots}

\pgfplotsset{every axis/.append style={
                    axis x line=middle,    % put the x axis in the middle
                    axis y line=middle,    % put the y axis in the middle
                    axis line style={<->,color=black}, % arrows on the axis
                    xlabel={$x$},          % default put x on x-axis
                    ylabel={$y$},          % default put y on y-axis
            }}
            
 \pgfplotsset{compat=1.14} % sometimes, everything falls apart without this
 
 %\usepgfplotslibrary{fillbetween} %for shading between curves

%%%%%% commands %%%%%%
\newcommand{\blue}[1]{{\color{blue}#1}}
\newcommand{\sol}[1]{~\\ \col{ {\em Solution: }#1} }
\newcommand{\col}[1]{{\color{blue}#1}} % I frequently use this to indicate where i"m going to be writing something
\newcommand{\solbf}[1]{~\\ \col{ {\textbf{Solution: }}#1} }
\newcommand{\al}[1]{\begin{align*} #1 \end{align*}} 
\newcommand{\itemI}[1]{\begin{itemize} #1 \end{itemize}} 
\def\ds{\displaystyle}

% exam heading; course name is first argument, exam information is second argument
\newcommand{\Heading}[2]{
    \fbox{
        {\Large\textbf{#1}}
        }
    \hfill
    \fbox{
        \begin{minipage}{0.2\textwidth}
            \begin{center}
            \textbf{#2}
            \end{center}
        \end{minipage}
    }
    \par\vspace{0.1in}}

% for being lazy while beautifying integrals
\newcommand{\dint}{\displaystyle\int}

% for being lazy while beautifying limits
\newcommand{\dlim}{\displaystyle\lim}

%%%%%% stuff %%%%%%

\usepackage{fancyhdr}

%\setcounter{page}{0} % first page is page 0

\parindent0pt % no indent

\fboxsep=0.25cm % little space in fbox border

%%%%%% BEGIN! %%%%%%%

\begin{document}

\pagestyle{fancy}

% record goals here to create sense of transparency

\fancyhead[c]{%
  \ifcase\value{page}
  % empty test for page = 0
  \or {}% 
  \or
  \or {\bf\color{ProcessBlue} F1: Lines \& Slope}% page=1
  \or {\bf\color{RoyalBlue} L1: Limits and Continuity}% page = 2
  \or {\bf\color{RoyalBlue} L2: Evaluating Limits Using Continuity and Forms}% page = 3
  \or {\bf\color{ForestGreen} D1: Meaning of a Derivative}% page = 4
  \or {\bf\color{ProcessBlue} F2: Exponential and Logarithmic Functions}% page = 5
  \or {\bf\color{ForestGreen} D3: Derivatives of Basic Functions}% page = 7
  \or {\bf\color{ForestGreen} D2: Compute Derivatives}% page = 6
  \or {\bf\color{RedOrange} A1: Derivatives and Graphing}% page = 8
  \or
  \else
  % Empty chead here!
  \fi
}

\thispagestyle{empty} % don't display that this is page 0

\Heading{MATH 1190/1210 Survey of Calculus I}{Knowledge Checkpoint 2\\Growth-Based\\Fall 2025}

\vspace{0.1in}

\textbf{Name} \hrulefill \, \begin{minipage}{0.16\textwidth}\begin{center}\textbf{UVA ID\\ (e.g. hdj4nd)} \end{center}\end{minipage}\hrulefill\\

\begin{minipage}{0.2\textwidth}\begin{center}\textbf{Course \&\\ Section Number} \end{center}\end{minipage} \hrulefill\,\,\textbf{Date} \hrulefill  \\

\hrulefill
\paragraph*{Course \& Section Numbers:}

{\small
\begin{center}
\begin{tabular}{| c | c | c | c | c |}
\hline
%
\begin{minipage}{0.12\textwidth}\begin{center}
1190-100\\
J. Rolf\\
(11:00 AM)
\end{center}\end{minipage}
&
\begin{minipage}{0.12\textwidth}\begin{center}
1190-200\\
J. Rolf\\
(9:30 AM)
\end{center}\end{minipage}
&
\begin{minipage}{0.12\textwidth}\begin{center}
1210-100\\
Mikhail\\
(8 AM)
\end{center}\end{minipage}
&
\begin{minipage}{0.12\textwidth}\begin{center}
1210-200\\
Raul\\
(9:30 AM)
\end{center}\end{minipage}
&
\begin{minipage}{0.12\textwidth}\begin{center}
1210-300\\
Sherman\\
(11 AM)
\end{center}\end{minipage}
\\
%
\hline
%
\begin{minipage}{0.12\textwidth}\begin{center}
1210-400\\
Sherman\\
(12:30 PM)
\end{center}\end{minipage}
&
\begin{minipage}{0.12\textwidth}\begin{center}
1210-500\\
Michael\\
(2 PM)
\end{center}\end{minipage}
&
\begin{minipage}{0.12\textwidth}\begin{center}
1210-600\\
J. Bowling\\
(3:30 PM)
\end{center}\end{minipage}
&
\begin{minipage}{0.12\textwidth}\begin{center}
1210-700\\
Josh\\
(9:30 AM)
\end{center}\end{minipage}
&
\begin{minipage}{0.12\textwidth}\begin{center}
1210-800\\
Mason\\
(2 PM)
\end{center}\end{minipage}
\\
%
\hline
%
\begin{minipage}{0.12\textwidth}\begin{center}
1210-900\\
Caelan\\
(3:30 PM)
\end{center}\end{minipage}
&
\begin{minipage}{0.12\textwidth}\begin{center}
1210-920\\
Caelan\\
(5 PM)
\end{center}\end{minipage}
&
\begin{minipage}{0.12\textwidth}\begin{center}
1210-940\\
Mitchell\\
(10 AM)
\end{center}\end{minipage}
&
\begin{minipage}{0.12\textwidth}\begin{center}
1210-950\\
Mitchell\\
(11 AM)
\end{center}\end{minipage}
&
\begin{minipage}{0.12\textwidth}\begin{center}
1210-960\\
Mitchell\\
(noon)
\end{center}\end{minipage}
\\
\hline
\end{tabular}
\end{center}

}

{%\small

\paragraph*{Instructions:}

\begin{itemize}
    
    \item The regular time allowed for completing this assessment is \fbox{$80$} minutes.
    %\item This assessment is out of 50 points.\\
    \item This assessment contains one single-part or multi-part problem for each of the following targets: {\bf\color{ProcessBlue}F1}, {\bf\color{RoyalBlue}L1}, {\bf\color{RoyalBlue}L2}, {\bf\color{ForestGreen}D1}, {\bf\color{ProcessBlue}F2}, {\bf\color{ForestGreen}D3}, {\bf\color{ForestGreen}D2}, {\bf\color{RedOrange}A1}.
    \item On each problem, you will receive a score of Exceeds Expectations (5), Proficient (4), Almost (3), Developing (2), Emerging (1), or Not Yet (0). {\bf Please do not interpret your score as a percentage.}
    %\item You will have opportunities to reassess on the targets appearing on this assessment after the next {\bf 2} assessments.\\
    \item You may NOT use calculators, dictionaries, notes, phones, or any other kind of aid.
    \item Please write neatly. What cannot be read, cannot be graded.
    \item Carefully work out each problem and clearly indicate your final answer to every problem.
    \item Throughout, give the most descriptive answer possible. \textbf{For instance, ``DNE" will not be accepted when ``$=\infty$" is applicable.}
    %\item Please provide the most specific answer possible. For instance, if a limit does not exist because it goes to $\infty$ or $-\infty$, then specify this.
    \item Please use the methods discussed up to this point in the course. For example, any work based on l'H\^opital's Rule will not receive credit.
    \item \textbf{All work must be shown unless otherwise specified.} Partial credit will not be given if appropriate work is not shown.
    \item In particular, {\bf any answer appearing with no justification will receive no credit regardless of whether the answer was correct}, unless otherwise specified.
    \item Please first respond to the prompts on which you feel most proficient, then respond to others later.
    \item This booklet is printed double-sided. It contains 12 pages and 8 numbered problems.
    \end{itemize}
    
    }

\vfill


%%%%%%%%%
\newpage

This page is left blank intentionally.\\

Use this page if you need additional space to write your solutions to any problem. \\

If you wish this page graded, you must refer to this page on \textbf{the original page} of the problem. Otherwise, your grader will NOT consider your work on this page.

\newpage%
%%%%%%%%%


\begin{enumerate}

%%%%%%%%%%%%%%%%%%%%%%%%%%%%%%%%%%%%%%%%%%%%%%%%%%%%%
%%%%%%%%%%% BEGIN PROBLEM F1 %%%%%%%%%%%%%%%%%%%%%%%%
%%%%%%%%%%%%%%%%%%%%%%%%%%%%%%%%%%%%%%%%%%%%%%%%%%%%%

\item The figure below shows a linear function representing a person's distance from their home (in miles) as a function of the amount of time they've spent on a trip (in hours).

\begin{center}
    \includegraphics[scale=0.6]{Fall 2025/KCs/KC2/KC2F1.png}
\end{center}

\begin{enumerate}
    \item How far away was the person from home when they started the trip?

    {\color{blue} \textbf{Solution:} At the start of the trip, they are 300 miles from home.    
    }

    \vfill

    \item Find the average rate of change of the person's distance from home from the beginning of the trip to 4 hours into the trip based on the graph. To receive full credit, please make sure to include appropriate units.

    {\color{blue} \textbf{Solution:} 
    \[\dfrac{y_2-y_1}{x_2-x_1}=\dfrac{500-300}{4-0}=\dfrac{200}{4}=50 mi/h\]
    }
    
    \vfill

    \item Use a sentence or two to explain the meaning of the average rate of change found in part (b) \textbf{in the context of this scenario}. 
    
    Note: Simply stating the ``slope" is the ``average rate of change" of the relevant quantities will not receive credit. You must explain the meaning of the slope using the details and quantities related to the context of this question.

    {\color{blue} \textbf{Solution:} In the first four hours, the person is 50 miles further from home for each hour that passes.
    
    Alternatively, we can say for each hour that passes; the person is 50 miles further from home for the first four hours.    
    }
    
    \vfill
\end{enumerate}

%%%%%%%%%%%%%%%%%%%%%%%%%%%%%%%%%%%%%%%%%%%%%%%%%%%%%
%%%%%%%%%%% END PROBLEM F1 %%%%%%%%%%%%%%%%%%%%%%%%%%
%%%%%%%%%%%%%%%%%%%%%%%%%%%%%%%%%%%%%%%%%%%%%%%%%%%%%

\newpage

%%%%%%%%%%%%%%%%%%%%%%%%%%%%%%%%%%%%%%%%%%%%%%%%%%%%%
%%%%%%%%%%% BEGIN PROBLEM L1 %%%%%%%%%%%%%%%%%%%%%%%%
%%%%%%%%%%%%%%%%%%%%%%%%%%%%%%%%%%%%%%%%%%%%%%%%%%%%%

\item For each of the following statements, clearly indicate whether it is true or false. If true, show graphical or verbal example(s) to demonstrate the statement. If false, fully explain why or show a counterexample.

\begin{enumerate}
    \item If $\lim\limits_{x\to -3}f(x)$ DNE, then both $\lim\limits_{x\to -3^-}f(x)$ and $\lim\limits_{x\to -3^+}f(x)$ must also DNE.
%%%%%%%%%%%%%%%%%%%%%%%%%%%%%%%%%%%%%%%%%%%%%%%%%%%%%
%%%%%%%%%%% L1 Part (a) Solution %%%%%%%%%%%%%%%%%%%%%%%%%%
%%%%%%%%%%%%%%%%%%%%%%%%%%%%%%%%%%%%%%%%%%%%%%%%%%%%%
\\
{
\color{blue}FALSE: We show by counterexample.
}
\begin{center}
\begin{tikzpicture}
  \begin{axis}[
    xlabel={$x$},
    ylabel={$f(x)$},
    xmin=-6, xmax=0,
    ymin=-1, ymax=5,
    grid=both,
    width=4in, height=3in,
    axis lines=middle,
    xtick={-3},
    ytick={1,2,3,4,5},
    clip=false
  ]
    % Left side: constant (jump)
    \addplot[blue, thick, domain=-6:-3, <-, >=stealth] {2};

    % Right side: vertical asymptote approaching +infinity
    \addplot[blue, thick, domain=-2.85:0, samples=100, ->, >=stealth] {1/(x+3)};

    % Open circle at jump
    \addplot[blue, fill=white, mark=o] coordinates {(-3,2)};

    % Vertical asymptote
    \addplot[dashed] coordinates {(-3,-1) (-3,6)};
    \node at (axis cs:-3,0.5) [anchor=south west] {$x=-3$};
  \end{axis}
\end{tikzpicture}
\end{center}

{
\color{blue}In this counterexample, we can see that the $\displaystyle \lim_{x \to -3}f(x)$ does not exist, but $\displaystyle \lim_{x \to -3^{+}}f(x)$ does exist even though $\displaystyle \lim_{x \to -3^{-}}f(x)$ does not. So, $\displaystyle \lim_{x \to 3} f(x)$ does not exist but $\displaystyle \lim_{x \to -3^{+}}f(x)$ does exist.
}

    \item If $f(x)$ has a removable discontinuity at $x=a$, then the continuity condition ``$f(a)$ is defined" is never true.
    \\
{
\color{blue}FALSE: We show by counter example.
}
\\
{
\color{Blue}
\begin{tikzpicture}
  \begin{axis}[
    xlabel={$x$},
    ylabel={$f(x)$},
    xmin=0, xmax=4,
    ymin=0, ymax=6,
    grid=both,
    width=3in, height=2.2in,
    samples=200,
    domain=0:4,
    axis lines=middle,
    xtick={2},
    ytick={4,5},
    clip=false
  ]
    % The main parabola (x^2)
    \addplot[blue, thick, domain=0:2.5, ->, >=stealth] {x^2};


    % Hole at (2,4)
    \addplot[blue, fill=white, mark=o] coordinates {(2,4)};
    % Defined point at (2,5)
    \addplot[blue, fill=blue, mark=*] coordinates {(2,5)};
  \end{axis}
\end{tikzpicture}

}
{
\color{blue}In this counterexample, we can see that the graph has a removable discontinuity at $x = 2$ ($\displaystyle \lim_{x \to 2^{-}}f(x) = \displaystyle \lim_{x \to 2^{+}}f(x) = 4$), and the function is defined at $x = 2$ ($f(2) = 5$). So, the function has a removable discontinuity at $x = 2$ and the statement ``$f(2)$ is defined" is true.
}
\end{enumerate}
%%%%%%%%%%%%%%%%%%%%%%%%%%%%%%%%%%%%%%%%%%%%%%%%%%%%%
%%%%%%%%%%% L1 Part (a) Solution %%%%%%%%%%%%%%%%%%%%%%%%%%
%%%%%%%%%%%%%%%%%%%%%%%%%%%%%%%%%%%%%%%%%%%%%%%%%%%%%

%%%%%%%%%%%%%%%%%%%%%%%%%%%%%%%%%%%%%%%%%%%%%%%%%%%%%
%%%%%%%%%%% END PROBLEM L1 %%%%%%%%%%%%%%%%%%%%%%%%%%
%%%%%%%%%%%%%%%%%%%%%%%%%%%%%%%%%%%%%%%%%%%%%%%%%%%%%

\newpage

%%%%%%%%%%%%%%%%%%%%%%%%%%%%%%%%%%%%%%%%%%%%%%%%%%%%%
%%%%%%%%%%% BEGIN PROBLEM L2 %%%%%%%%%%%%%%%%%%%%%%%%
%%%%%%%%%%%%%%%%%%%%%%%%%%%%%%%%%%%%%%%%%%%%%%%%%%%%%

\item Evaluate the following limits using a theorem about continuity or form.

If you use a theorem about continuity to evaluate a limit, to receive full credit, please provide sufficient explanation for why the theorem can be used. 

If you use form to evaluate a limit, to receive full credit, please list the form and show all steps. 

There's no need to simplify your final answer.

\begin{enumerate}
    \item $\lim\limits_{x\to\infty}\dfrac{x^5+4x^3-3\sqrt{x}}{4\log_2(x)-3x^5}$

    \vspace{.4in}

    {\color{blue}{As $x \to \infty$ the \underline{dominant term} in the numerator is $x^5$, and the \underline{dominant term} in the denominator is $-3x^5$.  Both numerator and denominator grow without bound, so this is the form $\frac{\infty}{-\infty}$.  The limit is the same as $\lim \limits_{x \to \infty} \frac{x^5}{-3x^5} = \boxed{-\frac13}$.}}

    \vspace{1in}

    \item $\lim\limits_{x\to 1}\dfrac{x^5+4x^3-3\sqrt{x}}{4\log_2(x)-3x^5}$

    \vspace{.4in}

    {\color{blue}{This function is built as a combination of polynomials, roots, and logarithms, so it is \underline{continuous} on its domain by the \underline{Main Continuity Theorem} of Lesson 3.}  The value $x=1$ does not require division by zero, or an even root of a negative number, or a logarithm of a nonpositive number, so it is in the domain of the function.  Thus by \underline{continuity} we have that the limit is the same as plugging in $x=1$, and the limit is $\frac{1^5 + 4 \cdot 1^3 - 3 \sqrt{1}}{4 \log_2 (1) - 3 \cdot 1^5} = \frac{2}{-3} = \boxed{-\frac23}$.}

    \vfill
\end{enumerate}

%%%%%%%%%%%%%%%%%%%%%%%%%%%%%%%%%%%%%%%%%%%%%%%%%%%%%
%%%%%%%%%%% END PROBLEM L2 %%%%%%%%%%%%%%%%%%%%%%%%%%
%%%%%%%%%%%%%%%%%%%%%%%%%%%%%%%%%%%%%%%%%%%%%%%%%%%%%

\newpage

%%%%%%%%%%%%%%%%%%%%%%%%%%%%%%%%%%%%%%%%%%%%%%%%%%%%%
%%%%%%%%%%% BEGIN PROBLEM D1 %%%%%%%%%%%%%%%%%%%%%%%%
%%%%%%%%%%%%%%%%%%%%%%%%%%%%%%%%%%%%%%%%%%%%%%%%%%%%%

\item Sam drives from home to school daily, where the commute takes approximately 20 minutes.

Below is a graph representing the relationship between the distance $s$ Sam has driven (in miles) and time $t$ (in minutes) during the commute.

In case it's unclear, there's no corner in the graph.

\begin{center}
        \includegraphics[scale=0.5]{Fall 2025/KCs/KC2/KC2D1.png}
    \end{center}

\begin{enumerate}
    \item Which one of i, ii, and iii best describes Sam's drive? Explain why.
    
    To receive full credit, point out at least two \textbf{graphical features related to derivative} that matches the description.

    \begin{enumerate}
        \item Sam drove on city roads where they stopped for a red light approximately every 5 minutes.
        \item Sam started driving, stopped at a stop sign in the neighborhood, then continued on a main road. They stopped again at a red light right before getting onto the highway. They drove at a constant speed on the highway for a 5-10 minutes before slowing down and eventually stopped as they arrived at the school.
        \item Sam started driving on a highway with a high speed for about 10 minutes, but then had to reduce the speed due to construction. Once they left the construction zone, they smoothly accelerated to a faster speed and drove at that speed for before arriving at the school and gradually come to a full stop.
    \end{enumerate}

    \textcolor{blue}{\textbf{Solution:} Choice ii best describes Sam's drive.\\
    The derivative $s'(t)$ represents the speed of Sam's driving in this scenario.\\
    At around 4 mins and around 10 mins, the graph appears to have horizontal tangent lines, indicating the derivative, or the speed, is $0$ at these times. This matches the two stops in choice ii.\\
    The slope of the graph (which represents the derivative) from about 4.5 mins to 9.5 mins is clearly less than the slope from about 11.5 mins to 18 mins. This matches the fact that Sam first drives on a main road and then a highway in choice ii, since the speed on a highway should be faster than the speed on a main road.}

    \item Rank the values of $s'(10)$, $s'(14)$, and $s'(19)$ from smallest to largest. No explanation required.

    \vspace{0.5cm}
    \textcolor{blue}{\textbf{Solution:} 
    \begin{center}
    \blue{$s'(10)<s'(19)<s'(14)$}
    \end{center}}
\end{enumerate}

%%%%%%%%%%%%%%%%%%%%%%%%%%%%%%%%%%%%%%%%%%%%%%%%%%%%%
%%%%%%%%%%% END PROBLEM D1 %%%%%%%%%%%%%%%%%%%%%%%%%%
%%%%%%%%%%%%%%%%%%%%%%%%%%%%%%%%%%%%%%%%%%%%%%%%%%%%%

\newpage

%%%%%%%%%%%%%%%%%%%%%%%%%%%%%%%%%%%%%%%%%%%%%%%%%%%%%
%%%%%%%%%%% BEGIN PROBLEM F2 %%%%%%%%%%%%%%%%%%%%%%%%
%%%%%%%%%%%%%%%%%%%%%%%%%%%%%%%%%%%%%%%%%%%%%%%%%%%%%

\item Respond to the following prompts. For this specific question, you are not required to compute any powers of constants or logarithms. For example, if your final answer includes $10^3$ or $\log_2(5)$, you can leave it as is.

\begin{enumerate}
    \item Simplify the expression below as much as possible so that the final result contains only positive exponents. You may expand exponential expressions and/or apply properties of exponents as needed. 

    \[\left(\dfrac{3^{-5/2}x}{\sqrt{3}x^{-3/5}}\right)^5\]
% -------------------------
% Problem (a)
% -------------------------
{
\color{blue}
\textbf{(a)} Simplify
\[
\left(\dfrac{3^{-5/2}x}{\sqrt{3}\,x^{-3/5}}\right)^5
\]

\begin{align*}
\text{First, rewrite }\sqrt{3}&=3^{1/2}\text{ and treat powers of }x\text{ separately.}\\
\frac{3^{-5/2}x}{\sqrt{3}\,x^{-3/5}}
&= \frac{3^{-5/2}\,x^{1}}{3^{1/2}\,x^{-3/5}} \\[6pt]
&= 3^{-5/2 - 1/2}\; x^{\,1 - (-3/5)} \qquad\text{(subtract exponents for same base)}\\[6pt]
&= 3^{-6/2}\; x^{1 + 3/5} \\[6pt]
&= 3^{-3}\; x^{8/5}.
\end{align*}

Now raise this to the 5th power using the power rule for exponents:
\[
\big(3^{-3}\; x^{8/5}\big)^5 = (3^{-3})^5 \cdot (x^{8/5})^5
= 3^{-15}\; x^{8}.
\]

Writing with only positive exponents yields
\[
\boxed{\dfrac{x^{8}}{3^{15}}}
\]
}
    \item Solve the following equation for $x$. 

    \[7^{2x-5}=10\]

% -------------------------
% Problem (b)
% -------------------------
\medskip
{
\color{blue}
\textbf{(b)} Solve for \(x\):
\[
7^{2x-5}=10.
\]

Take logarithms (base 7) of both sides:
\[
\log_7\!\big(7^{2x-5}\big)=\log_7(10).
\]

Use the logarithm cancellation rule:
\[
2x-5 = \log_7 10.
\]

Solve for \(x\):
\begin{align*}
2x-5 &= \log_7 10\\
2x &= 5 + \log_7 10.\\
x &= \frac{1}{2}\left(5 + \log_7 10.\right).
\end{align*}

\[
\boxed{\,x=\frac{5+\log_7(10)}{2}\,}
\]
}

\end{enumerate}


%%%%%%%%%%%%%%%%%%%%%%%%%%%%%%%%%%%%%%%%%%%%%%%%%%%%%
%%%%%%%%%%% END PROBLEM F2 %%%%%%%%%%%%%%%%%%%%%%%%%%
%%%%%%%%%%%%%%%%%%%%%%%%%%%%%%%%%%%%%%%%%%%%%%%%%%%%%

%\newpage

%%%%%%%%%%%%%%%%%%%%%%%%%%%%%%%%%%%%%%%%%%%%%%%%%%%%%
%%%%%%%%%%% BEGIN PROBLEM D3 %%%%%%%%%%%%%%%%%%%%%%%%
%%%%%%%%%%%%%%%%%%%%%%%%%%%%%%%%%%%%%%%%%%%%%%%%%%%%%

\item Write down the derivatives of the following functions. No work required. No need to simplify.

If you need to write down your work, please do so at the bottom half of this page, not in the answer space.

\begin{center}
        \begin{tabular}{l c l l c l}
            %
            $\displaystyle\frac{d}{dx}\left[ 9x^{6} \right]$ & $=$ & {\color{blue} $9 \cdot 6 x^5$ }&
            %
            $\displaystyle\frac{d}{dx}\left[ \dfrac{1}{\sqrt[5]{x^3}} \right]$ & $=$ &${\color{blue} -\displaystyle\frac{3}{5} x^{-\frac{8}{5}}}$\\[0.3in]
            %%%%%
            $\displaystyle\frac{d}{dx}\left[ \left(\dfrac{3}{4}\right)^x \right]$ & $=$ & ${\color{blue}  \left(\dfrac{3}{4}\right)^x \ln \left(\dfrac{3}{4}\right) } $& 
            %%%%%
            $\displaystyle\frac{d}{dx}\left[ \pi^{e} \right]$ & $=$ &$ {\color{blue} 0 }$ \\[0.3in]
            %%%%%
            $\displaystyle\frac{d}{dx}\left[ x^5+3 \right]$ & $=$ & ${\color{blue} 5x^4 }$&
            %
            $\displaystyle\frac{d}{dx}\left[ \sqrt{x}-\ln(x) \right]$ & $=$ & ${\color{blue} \dfrac{1}{2} x^{-\frac{1}{2}} - \dfrac{1}{x}}$\\[0.3in]
            %%%%%
            $\displaystyle\frac{d}{dx}\left[ \dfrac{x^3+5x^2-x}{x^4} \right]$ & $=$ & ${\color{blue} -x^{-2} - 5\cdot2 x^{-3}+3x^{-4}}$ & 
            %
            $\displaystyle\frac{d}{dx}\left[ \log_{2}(x) \right]$ & $=$ & ${\color{blue} \dfrac{1}{x \ln 2}}$\\[0.3in]
            %%%%%
        \end{tabular}        
\end{center}

{
    \color{blue} \textbf{Explanations} (not required to receive full credit)
    TBA
        
}


%%%%%%%%%%%%%%%%%%%%%%%%%%%%%%%%%%%%%%%%%%%%%%%%%%%%%
%%%%%%%%%%% END PROBLEM D3 %%%%%%%%%%%%%%%%%%%%%%%%%%
%%%%%%%%%%%%%%%%%%%%%%%%%%%%%%%%%%%%%%%%%%%%%%%%%%%%%

\newpage

%%%%%%%%%%%%%%%%%%%%%%%%%%%%%%%%%%%%%%%%%%%%%%%%%%%%%
%%%%%%%%%%% BEGIN PROBLEM D2 %%%%%%%%%%%%%%%%%%%%%%%%
%%%%%%%%%%%%%%%%%%%%%%%%%%%%%%%%%%%%%%%%%%%%%%%%%%%%%

\item Find the derivative of the following expressions. 

A correct final answer receives full credit. Nevertheless, you are encouraged to write down the steps so that you will be able to receive partial credit if you final answer is not fully correct. There's no need to simplify your final answer.
\begin{enumerate}
\item $\dfrac{\ln(x)+x^2}{x^3-2x}$\\
\solbf{ 

$f(x)=\dfrac{\ln(x)+x^2}{x^3-2x}$.
We see two functions that are divided. Therefore the quotient rule is needed.
\itemI{
\item  $\ds hi(x)=\ln(x)+x^2 \Rightarrow  hi^{\prime}(x)=\frac{1}{x}+2x$ 
\item  $lo(x)=x^3-2x \Rightarrow lo^{\prime}(x)=3x^2-2$
}

\al{
\frac{d}{dx}\left[ \dfrac{\ln(x)+x^2}{x^3-2x} \right]&=\frac{hi^{\prime}(x)\cdot lo(x)-hi(x)\cdot lo^{\prime}(x)}{lo^2(x)}\\
&=\frac{\left(\dfrac{1}{x}+2x \right)\left(x^3-2x \right)-\left(\ln(x)+x^2 \right) \left(3x^2-2 \right)}{\left(x^3-2x \right)^2}
}
}

\item $4\sqrt{x}\left(e^2+3x^4\right)^2$\\
\solbf{We see two functions that are multiplied together. Therefore the product rule is needed.
\itemI{
\item $f(x)=4\sqrt{x}=4x^{1/2}$
\item $g(x)=(e^2+3x^4)^2$
}
and $$4\sqrt{x}\left(e^2+3x^4\right)^2=f(x)g(x)$$

In order to implement the product rule, we are going to need $f^{\prime}$ and $g^{\prime}$
\itemI{
\item We use the power rule to find $f^{\prime}$
$$\ds f^{\prime}(x)=\frac{d}{dx}\left[4x^{1/2}\right]=4\frac{d}{dx}\left[x^{1/2}\right]=4\cdot \frac{1}{2}x^{-1/2} =2x^{-1/2} $$
\item To find $g^{\prime}$, we notice that we have a function-inside-a function. This will require the Chain Rule.
\al{
g^{\prime}(x)&=\underbrace{2\left( e^2+3x^4\right)^1}_{outer'}\cdot \underbrace{\left( 0+3\cdot 4x^3 \right)}_{inner'} \text{(since $e$ is a constant, $e^2$ is also a constant and thus $\dfrac{d}{dx}[e^2]=0$)}\\
&=2(e^2+3x^4)(12x^3)
}
}
We put all of this together.
\al{
\frac{d}{dx}\left[4\sqrt{x}\left(e^2+3x^4\right)^2\right]& =\frac{d}{dx}\left[ f(x)g(x)\right]\\
&=f^{\prime}(x)\cdot g(x)+f(x)\cdot g^{\prime}(x)\\
&=\left( 2x^{-1/2}\right)\cdot (e^2+3x^4)^2+\left(4x^{1/2}\right) \cdot 2(e^2+3x^4)(12x^3)
}
}

\end{enumerate}


%split part (a) and part (b)

%%%%%%%%%%%%%%%%%%%%%%%%%%%%%%%%%%%%%%%%%%%%%%%%%%%%%
%%%%%%%%%%% END PROBLEM D2 %%%%%%%%%%%%%%%%%%%%%%%%%%
%%%%%%%%%%%%%%%%%%%%%%%%%%%%%%%%%%%%%%%%%%%%%%%%%%%%%

\newpage

%%%%%%%%%%%%%%%%%%%%%%%%%%%%%%%%%%%%%%%%%%%%%%%%%%%%%
%%%%%%%%%%% BEGIN PROBLEM A1 %%%%%%%%%%%%%%%%%%%%%%%%
%%%%%%%%%%%%%%%%%%%%%%%%%%%%%%%%%%%%%%%%%%%%%%%%%%%%%

\item Consider the following function $f(x)$.

\[f(x)=x^3-6x^2+5\]

Respond to the following prompts. Make sure to show all work.

\begin{enumerate}
    \item Find all critical numbers of $f(x)$.

    {\color{blue}\textbf{Solution:} $f'(x)=3x^2-12x$.\\
    If $f'(x)=0$, we get 
    \begin{align*}
        3x^2-12x=0\\
        x(3x-12)=0\\
        x=0\text{ or } 3x-12=0\\
        x=0\text{ or } 3x=12\\
        x=0, 4
    \end{align*}
    
    There's no $x$-value at which $f'(x)$ DNE, since $f'(x)$ is a polynomial and therefore is defined for all real numbers.
    
    Since both $x=0$ and $x=4$ are in the domain of $f(x)$, they are critical numbers of $f(x)$.}

    \vfill

    \item On what interval(s) is $f(x)$ concave up? On what intervals(s) is $f(x)$ concave down?

    {\color{blue}\textbf{Solution:} 
    \begin{align*}
        f''(x)=6x=12\\
        6x-12=0\\
        6x=12\\
        x=2\\
    \end{align*}
    So the intervals of interest are $(-\infty, 2)$ and $(2, \infty)$.\\    
    $f''(1)=6-12=-6<0$ and $f''(3)=18-12=6>0$.\\
    So $f(x)$ is concave down on $(-\infty, 2)$ and concave up on $(2, \infty)$.}

    \vfill
\end{enumerate}


%%%%%%%%%%%%%%%%%%%%%%%%%%%%%%%%%%%%%%%%%%%%%%%%%%%%%
%%%%%%%%%%% END PROBLEM A1 %%%%%%%%%%%%%%%%%%%%%%%%%%
%%%%%%%%%%%%%%%%%%%%%%%%%%%%%%%%%%%%%%%%%%%%%%%%%%%%%

\newpage

This page is left blank intentionally.\\

Use this page if you need additional space to write your solutions to any problem. \\

If you wish this page graded, you must refer to this page on \textbf{the original page} of the problem. Otherwise, your grader will NOT consider your work on this page.

\newpage

This page is left blank intentionally.\\

Use this page if you need additional space to write your solutions to any problem. \\

If you wish this page graded, you must refer to this page on \textbf{the original page} of the problem. Otherwise, your grader will NOT consider your work on this page.

\end{enumerate}

\end{document}

Important idea 1: Clear indication of whether statement is T or F
For a) clearly shows/explains counterexample with the three properties
Important idea 2:	Limit DNE
Important idea 3: At least one one-sided limit does exist 
For b) Clearly shows/explains counterexample with two properties
Important idea 4: Has  a removable discontinuity
Important idea 5:  Function is defined at the point of discontinuity.
